\usection{Validación del Lenguaje y Entorno}

Para garantizar que el lenguaje de programación visual basado en flujos de datos cumpliera con los objetivos de investigación planteados, se diseñó y ejecutó un proceso de validación estructurado en dos etapas fundamentales. La primera etapa se enfocó en evaluar la efectividad pedagógica de la herramienta a través del juicio de un experto, mientras que la segunda etapa consistió en pruebas de usabilidad directas con el público objetivo.

\usubsubsection{Evaluación de la efectividad pedagógica en el contexto educativo}

Con el propósito de validar la pertinencia del sistema como instrumento para el fomento del pensamiento computacional, se llevó a cabo una sesión de evaluación cualitativa con la profesora Mary Acosta, docente con experiencia en el área del pensamiento lógico-matemático.

Durante esta sesión de validación, se le presentó detalladamente el entorno de desarrollo a la docente. Se procedió a explicar el paradigma de programación por flujos de datos (dataflow), la gramática visual de los nodos y las potencialidades del lenguaje para representar procesos lógicos y matemáticos concurrentes.

Tras interactuar con el sistema y analizar su estructura, la profesora Acosta emitió una valoración altamente positiva respecto a la aplicabilidad del proyecto. Expresó que la herramienta tiene un lugar claro y de gran utilidad dentro de un aula de clases regular, destacando que el enfoque visual permite a los estudiantes estructurar sus ideas y desarrollar el pensamiento lógico de forma progresiva e intuitiva. De manera particular, la docente resaltó el alto potencial de la aplicación para la enseñanza y el refuerzo de contenidos matemáticos. Al visualizar cómo los datos numéricos fluyen a través de los nodos de procesamiento, los estudiantes pueden tangibilizar conceptos matemáticos abstractos, facilitando un aprendizaje mucho más significativo y participativo en esta área.

\usubsubsection{Pruebas de usabilidad con el grupo objetivo}

La segunda etapa de validación se centró en evaluar la usabilidad y la experiencia de usuario (UX) del entorno de ejecución directamente con su población objetivo: niños y jóvenes de 8 a 16 años.

Se organizaron sesiones de prueba controladas donde se le solicitó a un grupo de niños interactuar con el editor visual mediante la realización de diferentes tareas de programación predefinidas. Durante estas dinámicas, se aplicó la técnica de observación directa para analizar cómo los usuarios se desenvolvían en el entorno: cómo arrastraban los nodos al lienzo, cómo trazaban las conexiones entre los puertos y cómo interpretaban la retroalimentación visual del sistema al ejecutar los programas, incluyendo la gestión de errores y la depuración de programas como elemento funcionalidad del uso del entorno.

A lo largo de las sesiones, se recopilaron los comentarios, dudas y observaciones espontáneas emitidas por los niños. El análisis de estas interacciones permitió verificar de forma concluyente que la interfaz gráfica diseñada es altamente compatible con las capacidades motoras y cognitivas de los niños. El entorno demostró ser intuitivo, confirmando que la representación de un programa como un grafo de conexiones lógicas resulta asimilable y atractiva para los estudiantes, mitigando las frustraciones típicas asociadas a la sintaxis estricta de los lenguajes de programación basados en texto.

\usubsubsection{Resultados de la validación}

En conjunto, los resultados de ambas etapas de validación confirmaron el éxito del desarrollo. Desde la perspectiva experta, el sistema se consolida como un recurso didáctico de alto valor para las aulas de clase, especialmente en el área de matemáticas y lógica; y desde la perspectiva del usuario final, la herramienta cumple exitosamente con los requerimientos no funcionales de usabilidad, garantizando que el público infantil y juvenil pueda utilizarla de forma adecuada para iniciarse en el pensamiento computacional.